\documentclass[11pt,a4paper]{article}

\usepackage[utf8]{inputenc}
\usepackage[T1]{fontenc}
\usepackage[brazilian]{babel}
\usepackage{geometry}
\usepackage{xcolor}
\usepackage{tabularx}
\usepackage{longtable}
\usepackage{booktabs}
\usepackage{array}
\usepackage{enumitem}
\usepackage{hyperref}
\usepackage{titlesec}
\usepackage{fancyhdr}
\usepackage{lastpage}

\geometry{margin=2.2cm}
\hypersetup{
  colorlinks=true,
  linkcolor=black,
  urlcolor=blue!60!black
}

\definecolor{artemisBlue}{HTML}{17324D}
\definecolor{artemisGreen}{HTML}{217A5B}
\definecolor{artemisGold}{HTML}{B8872B}
\definecolor{artemisRed}{HTML}{9B2D20}
\definecolor{lightGray}{HTML}{F3F5F7}

\titleformat{\section}
  {\Large\bfseries\color{artemisBlue}}
  {\thesection}{0.8em}{}

\titleformat{\subsection}
  {\large\bfseries\color{artemisBlue}}
  {\thesubsection}{0.8em}{}

\titleformat{\subsubsection}
  {\normalsize\bfseries\color{artemisBlue}}
  {\thesubsubsection}{0.8em}{}

\pagestyle{fancy}
\fancyhf{}
\setlength{\headheight}{14pt}
\lhead{\textbf{Artemis Bot}}
\rhead{Auditoria e proposta de conclusao}
\cfoot{\thepage\ de \pageref{LastPage}}

\emergencystretch=2em
\sloppy

\setlist[itemize]{leftmargin=1.2em}
\setlist[enumerate]{leftmargin=1.4em}
\renewcommand{\arraystretch}{1.25}

\newcommand{\statusok}{\textcolor{artemisGreen}{Implementado}}
\newcommand{\statuspartial}{\textcolor{artemisGold}{Parcial}}
\newcommand{\statusrisk}{\textcolor{artemisRed}{Critico}}
\newcommand{\statusmissing}{\textcolor{artemisRed}{Pendente}}

\begin{document}

\begin{titlepage}
  \centering
  \vspace*{2.0cm}
  {\Huge\bfseries\color{artemisBlue} Artemis Bot\par}
  \vspace{0.4cm}
  {\Large Relatorio executivo, auditoria comparativa e proposta de conclusao\par}
  \vspace{1.2cm}
  {\large Assistente de IA para WhatsApp com integrações de CRM, agenda e pagamento\par}
  \vfill
  \begin{tabular}{rl}
    \textbf{Cliente:} & Confluence Treinamento \\
    \textbf{Base analisada:} & Repositorio local Artemis-Bot + planejamento em PDF \\
    \textbf{Data:} & 24/04/2026 \\
    \textbf{Status:} & Desenvolvimento avancado, pendente de endurecimento tecnico \\
  \end{tabular}
  \vfill
  {\small Documento preparado para avaliacao comercial e tecnica. Valores e prazos sao estimativas e devem ser confirmados apos fechamento formal de escopo.}
\end{titlepage}

\tableofcontents
\newpage

\section{Mensagem executiva ao cliente}

O Artemis Bot ja ultrapassou a fase de ideia. O que existe hoje nao e apenas um chatbot simples: e uma aplicacao TypeScript sob medida, integrada ao WhatsApp oficial, banco de dados, IA generativa, Google Calendar, Asaas e formulario Respondi. Isso significa que o projeto ja tem uma base real de produto, com automacao de atendimento, qualificacao comercial, agendamento e cobranca.

O ponto central da avaliacao e este: o projeto foi iniciado com apoio de IA e ja possui uma quantidade relevante de implementacao, mas ainda precisa passar por uma etapa profissional de conclusao. Essa etapa nao e simplesmente ``continuar escrevendo codigo''. O trabalho correto agora e transformar uma base funcional em uma entrega confiavel, segura, documentada e defensavel perante o cliente final.

Como CEO ou responsavel pela entrega, eu apresentaria o Artemis da seguinte forma:

\begin{quote}
  O projeto Artemis ja tem o nucleo do produto implementado. Ele conversa pelo WhatsApp, entende o estagio do lead, salva historico, qualifica o aluno, apresenta programas, trata objecoes, agenda no calendario e gera cobrancas. A etapa restante e a mais importante para uma entrega profissional: seguranca, testes, validacao de fluxos, estabilizacao de integracoes, documentacao operacional e preparacao para producao.
\end{quote}

O valor agregado nao esta apenas no bot responder mensagens. O valor esta em reduzir trabalho manual, responder leads com consistencia, preservar contexto, encaminhar o cliente pelo funil comercial e automatizar pontos que normalmente dependem de uma pessoa: cadastro, agenda, pagamento e confirmacao.

\section{Resumo de valor entregue}

\subsection{O que o cliente recebe}

\begin{itemize}
  \item Um assistente de WhatsApp com IA personalizado para a Confluence Treinamento.
  \item Uma arquitetura propria em Node.js e TypeScript, sem depender de construtores visuais limitados.
  \item Memoria persistente de conversas em PostgreSQL.
  \item Fluxo comercial por estados: saudacao, qualificacao, apresentacao, objecao, fechamento e humano.
  \item Integracao com Meta WhatsApp Business API.
  \item Integracao com Google Gemini para respostas e extracao de dados.
  \item Integracao com Google Calendar para agenda.
  \item Integracao com Asaas para cliente, cobranca, assinatura e cancelamento.
  \item Integracao com Respondi para formulario de matricula.
  \item Regras de LGPD: consentimento, exclusao de dados e retencao de leads inativos.
  \item Configuracao de persona, programas e mensagens sem alterar codigo.
\end{itemize}

\subsection{Por que isso tem valor comercial}

Um atendimento humano pode perder leads por demora, inconsistência ou falta de follow-up. O Artemis tenta resolver esse gargalo com disponibilidade continua e processo padronizado. Alem disso, ao integrar agenda e pagamento, o sistema reduz a distancia entre ``tenho interesse'' e ``matricula iniciada''.

Na pratica, o ativo entregue e uma automacao comercial com IA, nao apenas uma conversa automatica. Isso justifica tratar o projeto como software sob medida, com fase de estabilizacao e manutencao, e nao como configuracao simples de chatbot.

\section{Estimativa comercial de conclusao}

\subsection{Premissas}

A estimativa abaixo considera que o cliente ja trouxe uma base iniciada e que o trabalho contratado agora seria concluir, estabilizar e preparar a entrega para uso real. Nao inclui custos de terceiros, como Meta, Asaas, Google, VPS, dominio, certificados, banco gerenciado ou mensalidade de APIs.

Referencias publicas recentes de mercado indicam ampla variacao de custo para desenvolvimento sob medida e horas tecnicas senior no Brasil. Foram usadas apenas como baliza externa, nao como tabela obrigatoria: ha fontes citando hora de desenvolvimento na faixa de R\$ 120--180/h para software sob medida, faixa ampla de freelancer autonomo entre R\$ 68--243/h e projetos de software variando muito conforme risco e complexidade.\footnote{\url{https://www.aegisai.com.br/blog/preco-desenvolvimento-software-sob-medida-2026}}\footnote{\url{https://www.freelasemcrise.com.br/quanto-cobra/desenvolvedor-software/sao-paulo}}\footnote{\url{https://www.neryx.com.br/blog/custo-desenvolvimento-software-brasil-2026/}}

\subsection{Escopo recomendado para fechar o projeto}

\begin{longtable}{p{0.28\textwidth}p{0.18\textwidth}p{0.44\textwidth}}
\toprule
\textbf{Frente} & \textbf{Esforco estimado} & \textbf{Entrega} \\
\midrule
Auditoria tecnica e organizacao & 12--20 h & Revisao da base, definicao de riscos, ajuste de README, envs e runbook. \\
Seguranca e producao & 24--40 h & Validacao webhook Asaas, remocao de logs sensiveis, graceful shutdown, Calendar ID por ambiente. \\
Banco e infraestrutura de app & 16--28 h & Prisma Singleton, migrations, historico correto, ajustes de deploy. \\
Testes e QA & 32--56 h & Testes unitarios de FSM, normalizador, pagamentos, prompts e checklist end-to-end. \\
Refinamento de funil & 24--44 h & Ajuste de protocolo teen, PRM, midias nao suportadas, fechamento e handoff. \\
Deploy assistido e transferencia & 16--28 h & Subida controlada, validacao sandbox, documentacao operacional e treinamento. \\
\bottomrule
\end{longtable}

\subsection{Proposta de investimento}

Para uma entrega profissional, eu recomendaria apresentar tres opcoes:

\begin{longtable}{p{0.20\textwidth}p{0.18\textwidth}p{0.22\textwidth}p{0.30\textwidth}}
\toprule
\textbf{Plano} & \textbf{Prazo} & \textbf{Investimento} & \textbf{Quando usar} \\
\midrule
Essencial & 2--3 semanas & R\$ 18.000--26.000 & Corrigir riscos criticos, documentar e deixar rodando com QA manual limitado. \\
Profissional & 4--6 semanas & R\$ 32.000--48.000 & Recomendado. Inclui seguranca, testes, QA, deploy assistido e documentacao de operacao. \\
Produto & 8--10 semanas & R\$ 58.000--85.000 & Inclui melhorias de produto, painel/observabilidade inicial, status de matricula e preparacao para evolucao SaaS. \\
\bottomrule
\end{longtable}

\subsection{Preco recomendado para esta oportunidade}

Para este caso especifico, considerando que ja existe uma base funcional e que o cliente quer concluir com seguranca, minha recomendacao comercial seria:

\begin{center}
\fcolorbox{artemisBlue}{lightGray}{
  \begin{minipage}{0.86\textwidth}
    \centering
    \vspace{0.2cm}
    {\Large\bfseries Proposta recomendada: R\$ 38.000}\\
    \vspace{0.2cm}
    Prazo estimado: 5 semanas\\
    Escopo: estabilizacao tecnica, seguranca, testes essenciais, QA completo do funil, documentacao e deploy assistido.\\
    \vspace{0.2cm}
  \end{minipage}
}
\end{center}

Manutencao opcional apos entrega:

\begin{itemize}
  \item R\$ 2.500--3.500/mês para suporte leve, ajustes de prompt e acompanhamento.
  \item R\$ 4.500--6.500/mês para suporte com evolucoes mensais, monitoramento e melhorias continuas.
\end{itemize}

\subsection{Como defender o preco}

O argumento nao deve ser ``vou terminar um bot''. O argumento correto e:

\begin{quote}
  Vamos transformar um prototipo avancado em um sistema operacional de vendas e atendimento, com integracoes reais de calendario, pagamento, dados e WhatsApp oficial. O trabalho inclui seguranca, testes, estabilidade e transferencia para que a Confluence consiga usar isso como ativo de negocio, nao como experimento.
\end{quote}

\section{Veredito tecnico resumido}

O projeto real esta muito proximo do planejamento em arquitetura geral, stack, integracoes principais e intencao de produto. A base tem os 16 arquivos TypeScript citados no planejamento e soma 2.728 linhas de TypeScript em \texttt{src/}, alinhado com o documento original.

\subsection{Pontos fortes}

\begin{itemize}
  \item Estrutura modular clara: controllers, services, flow, types e utils.
  \item FSM de conversa implementada.
  \item Debounce por usuario no webhook principal.
  \item Prompt dinamico por estado em XML.
  \item Gemini com tool calling para agenda e pagamentos.
  \item Asaas para cliente, cobranca, assinatura e cancelamento.
  \item Google Calendar para disponibilidade, criacao, busca e cancelamento.
  \item Respondi com validacao por token e normalizacao de telefone.
  \item Banco PostgreSQL via Prisma.
  \item Preocupacao real com LGPD.
\end{itemize}

\subsection{Riscos antes de producao}

\begin{itemize}
  \item Webhook Asaas aceita eventos sem assinatura, secret ou verificacao de origem.
  \item Existem multiplas instancias diretas de \texttt{PrismaClient}.
  \item Nao existe suite de testes.
  \item Nao existem migrations Prisma versionadas.
  \item Nao existe \texttt{.env.example}, configuracao PM2 ou runbook de deploy completo no repositorio.
  \item README esta desatualizado e cita OpenAI/variaveis antigas.
  \item Calendar ID esta hardcoded para um email pessoal/teste.
  \item Nao ha graceful shutdown.
  \item Parte do comportamento critico depende de prompt, nao de validacao backend.
\end{itemize}

\section{Inventario real do projeto}

\subsection{Arquivos encontrados}

\begin{longtable}{p{0.42\textwidth}p{0.48\textwidth}}
\toprule
\textbf{Arquivo} & \textbf{Funcao no projeto} \\
\midrule
\texttt{package.json} & Define scripts, dependencias e entrada de producao. Contem Express, Prisma, Gemini, Google APIs, Axios e utilitarios. \\
\texttt{package-lock.json} & Trava versoes instaladas das dependencias npm. \\
\texttt{tsconfig.json} & Configura compilacao TypeScript para CommonJS, target ES2020 e saida em \texttt{dist/}. \\
\texttt{prisma.config.ts} & Configuracao Prisma apontando para \texttt{prisma/schema.prisma} e \texttt{DATABASE\_URL}. \\
\texttt{README.md} & Documentacao inicial, porem desatualizada em variaveis e provedor de IA. \\
\texttt{prisma/schema.prisma} & Define modelos \texttt{User} e \texttt{ChatHistory}. \\
\texttt{config/persona.json} & Persona, tom, restricoes, objecoes, contratos e links de protocolo. \\
\texttt{config/programs.json} & Programas vendidos, precos, duracao e textos informativos. \\
\texttt{config/settings.json} & Mensagens e numeros de suporte/handoff. \\
\texttt{src/index.ts} & Inicializa Express, registra rotas e cron LGPD. \\
\texttt{src/controllers/WebhookController.ts} & Webhook principal WhatsApp, debounce, processamento background, handoff e envio de respostas. \\
\texttt{src/controllers/AsaasWebhookController.ts} & Recebe eventos Asaas e envia confirmacao de pagamento ao usuario. \\
\texttt{src/controllers/RespondiController.ts} & Recebe formulario Respondi, valida token, extrai campos e faz upsert do usuario. \\
\texttt{src/flow/StateResolver.ts} & Resolve estado conversacional e tipos de transferencia humana. \\
\texttt{src/services/AIService.ts} & Integra Gemini, extrai perfil, gera resposta e executa ferramentas de agenda/pagamento. \\
\texttt{src/services/AsaasService.ts} & Busca/cria clientes, gera cobrancas/assinaturas e cancela pendencias no Asaas. \\
\texttt{src/services/CalendarService.ts} & Verifica disponibilidade, cria, busca e cancela eventos no Google Calendar. \\
\texttt{src/services/ConfigLoader.ts} & Carrega arquivos JSON de configuracao. \\
\texttt{src/services/PromptBuilder.ts} & Monta prompt XML por estado da conversa. \\
\texttt{src/services/StateService.ts} & Gerencia usuario, historico, perfil, estado e exclusao LGPD. \\
\texttt{src/services/WhatsAppService.ts} & Envia mensagens pela Meta Graph API. \\
\texttt{src/types/user.ts} & Define estados de conversa e tipos de usuario. \\
\texttt{src/types/config.ts} & Define tipos TypeScript para persona, programas e settings. \\
\texttt{src/utils/prisma.ts} & Singleton Prisma para acesso ao banco. \\
\texttt{src/utils/phoneNormalizer.ts} & Normaliza telefones brasileiros para formato canonico. \\
\bottomrule
\end{longtable}

\subsection{Contagem local}

\begin{tabularx}{\textwidth}{Xr}
\toprule
\textbf{Grupo} & \textbf{Linhas} \\
\midrule
TypeScript em \texttt{src/} & 2.728 \\
Config JSON & 119 \\
Prisma schema & 50 \\
README & 51 \\
Total auditado localmente & 3.059 \\
\bottomrule
\end{tabularx}

\section{Matriz comparativa}

\begin{longtable}{p{0.25\textwidth}p{0.16\textwidth}p{0.49\textwidth}}
\toprule
\textbf{Area} & \textbf{Status} & \textbf{Leitura tecnica} \\
\midrule
TypeScript/Node.js & \statusok & Stack real bate com o planejamento. \\
Express HTTP API & \statusok & Rotas principais existem. \\
PostgreSQL + Prisma & \statuspartial & Schema existe, mas sem migrations versionadas. \\
User + ChatHistory & \statusok & Campos batem com o planejamento. \\
FSM com 6 estados & \statusok & Estado e resolucao implementados em TypeScript. \\
Debounce 5s & \statusok & Buffer por usuario existe. \\
Resposta 200 imediata Meta & \statusok & ACK acontece antes do processamento pesado. \\
PromptBuilder XML & \statusok & Ponto forte da arquitetura. \\
Gemini principal & \statusok & Usa \texttt{gemini-3.1-pro-preview}. \\
Gemini extracao JSON & \statusok & Usa \texttt{gemini-2.5-flash}. \\
Tool calling 6 ferramentas & \statusok & 4 ferramentas Calendar + 2 Asaas. \\
Filtro de raciocinio & \statuspartial & Existe heuristica, mas nao e garantia formal. \\
Google Calendar & \statuspartial & Funcional, mas Calendar ID hardcoded. \\
Asaas & \statuspartial & Cobranca existe; webhook sem validacao e risco critico. \\
Respondi & \statusok & Valida token e faz upsert. \\
LGPD exclusao & \statusok & Delecao com cascade implementada. \\
LGPD retencao 30 dias & \statuspartial & Cron existe, mas usa Prisma separado e \texttt{setInterval}. \\
Graceful shutdown & \statusmissing & Nao implementado. \\
Testes & \statusmissing & Sem arquivos de teste. \\
QA end-to-end & \statusmissing & Sem cenarios automatizados ou evidencias. \\
SaaS multi-tenant & \statusmissing & Roadmap futuro, nao implementado. \\
\bottomrule
\end{longtable}

\section{Comparativo por area funcional}

\subsection{Arquitetura principal}

O planejamento previa uma aplicacao em camadas com FSM, debounce, PromptBuilder e ConfigLoader. Isso esta implementado. O codigo separa controllers, services, flow, types e utils, o que e positivo para manutencao.

O ponto de atencao e que a serializacao por usuario ainda e limitada ao debounce. Se uma conversa gerar processamento longo e o usuario mandar outra rajada, pode haver concorrencia entre processamentos do mesmo numero. Para producao, uma fila por usuario seria mais segura.

\subsection{Inteligencia artificial}

O projeto usa Gemini em duas funcoes: resposta principal e extracao de perfil. A resposta principal recebe prompt dinamico por estado e pode chamar ferramentas. A extracao tenta preencher nome, idade, objetivo e alvo da matricula.

O desenho e bom, mas regras criticas ainda estao concentradas no prompt. Em sistemas com pagamento, agenda e LGPD, o backend deve validar regras essenciais antes de executar a acao.

\subsection{Funil comercial}

O funil existe e cobre saudacao, qualificacao, apresentacao, objecoes, fechamento e transferencia humana. A maior divergencia e que a ordem atual de qualificacao coleta nome, alvo da matricula, idade e objetivo; o planejamento falava em objetivo e idade primeiro.

Tambem existe diferenca na regra teen: o documento fala 14--16 em alguns trechos, enquanto o prompt usa 12--17 em outros. Essa regra precisa ser unificada antes da entrega.

\subsection{Asaas}

O Asaas esta bem encaminhado para criacao de cliente, cobranca, assinatura, link e cancelamento. O maior problema e o webhook: hoje qualquer origem poderia simular pagamento confirmado se souber o endpoint e o formato do payload.

Esse e o principal item de seguranca antes de producao.

\subsection{Google Calendar}

As quatro operacoes principais existem. O sistema verifica disponibilidade antes de criar evento e suporta recorrencia semanal.

Problemas atuais:
\begin{itemize}
  \item Calendar ID hardcoded.
  \item Credencial esperada em arquivo local nao documentado adequadamente.
  \item Regras de horario comercial ficam no prompt, nao no backend.
\end{itemize}

\subsection{Respondi}

O Respondi esta implementado de forma funcional. O controller valida token, extrai respostas por titulo, normaliza telefone e atualiza/cria usuario.

O principal ajuste e reduzir logs com dados pessoais. Payload bruto com CPF, email e endereco nao deve ir para log de producao.

\subsection{LGPD}

Ha uma preocupacao real com LGPD: consentimento, exclusao por solicitacao, cascade delete e retencao de leads inativos.

Ainda faltam:
\begin{itemize}
  \item Fluxo deterministico para acesso/correcao de dados.
  \item Politica de log sem dados pessoais.
  \item Retencao clara para clientes convertidos.
\end{itemize}

\section{O que ja foi feito}

\begin{itemize}
  \item Servidor Express com rotas principais.
  \item Webhook WhatsApp com ACK imediato.
  \item Debounce por usuario de 5 segundos.
  \item Persistencia de usuario e historico.
  \item Normalizacao de telefone brasileiro.
  \item FSM conversacional.
  \item Extracao de perfil com Gemini.
  \item Resposta principal com Gemini.
  \item Prompt dinamico por estado.
  \item Tool calling com Calendar e Asaas.
  \item Envio de WhatsApp via Meta Graph API.
  \item Intake Respondi.
  \item Cadastro, assinatura, cobranca e cancelamento no Asaas.
  \item Confirmacao de pagamento via webhook.
  \item Integracao Google Calendar.
  \item Exclusao LGPD.
  \item Retencao de leads inativos por 30 dias.
  \item Configuracao por JSON.
\end{itemize}

\section{O que falta fazer}

\subsection{Prioridade critica}

\begin{enumerate}
  \item Validar webhook Asaas.
  \item Remover instancias diretas de \texttt{new PrismaClient()}.
  \item Adicionar graceful shutdown.
  \item Corrigir busca das ultimas 20 mensagens.
  \item Criar testes unitarios para regras centrais.
  \item Criar \texttt{.env.example}.
  \item Criar migrations Prisma versionadas.
\end{enumerate}

\subsection{Prioridade alta}

\begin{enumerate}
  \item Mover Calendar ID para variavel de ambiente.
  \item Validar horario comercial no backend.
  \item Adicionar idempotencia no webhook Asaas.
  \item Remover logs de payload bruto com dados pessoais.
  \item Atualizar README.
  \item Criar checklist QA end-to-end.
  \item Criar script/config PM2.
\end{enumerate}

\subsection{Prioridade media}

\begin{enumerate}
  \item Transformar regras criticas de prompt em validacoes backend.
  \item Adicionar status de matricula/pagamento no banco.
  \item Melhorar roteamento deterministico para PRM.
  \item Tipar melhor sessoes e perfis.
  \item Remover codigo morto.
  \item Criar logs estruturados.
\end{enumerate}

\section{Plano de entrega recomendado}

\begin{longtable}{p{0.18\textwidth}p{0.24\textwidth}p{0.48\textwidth}}
\toprule
\textbf{Semana} & \textbf{Foco} & \textbf{Entregas} \\
\midrule
1 & Organizacao e riscos & README, envs, Prisma Singleton, historico correto, plano QA. \\
2 & Seguranca & Webhook Asaas protegido, logs sensiveis removidos, Calendar env. \\
3 & Testes essenciais & FSM, normalizador, PromptBuilder, calculos Asaas, Respondi. \\
4 & Funil e integracoes & Ajuste teen/PRM, midias nao suportadas, QA sandbox Calendar/Asaas. \\
5 & Deploy e transferencia & Build limpo, PM2, runbook, checklist final, treinamento e aceite. \\
\bottomrule
\end{longtable}

\section{Criterios de aceite}

\begin{itemize}
  \item Build TypeScript limpo.
  \item Migrations aplicadas por fluxo controlado.
  \item Webhook Asaas protegido e idempotente.
  \item Logs sem dados pessoais sensiveis.
  \item QA manual documentado.
  \item Testes unitarios cobrindo regras centrais.
  \item Graceful shutdown.
  \item Configuracao de producao sem Calendar ID hardcoded.
  \item README/runbook coerente com o codigo real.
  \item Fluxos principais validados em sandbox.
\end{itemize}

\section{Conclusao}

O Artemis Bot tem uma base valiosa e funcional. Ele cobre grande parte do planejamento original e mostra que o produto pode virar uma automacao comercial real para a Confluence. Ao mesmo tempo, ainda nao deve ser vendido como producao madura sem a etapa de endurecimento tecnico.

A proposta correta e posicionar a entrega como conclusao profissional de um sistema ja iniciado: corrigir riscos, proteger integracoes financeiras, documentar operacao, testar fluxos e preparar deploy. Com esse enquadramento, a proposta recomendada de R\$ 38.000 por aproximadamente 5 semanas e defensavel, desde que o escopo seja fechado e os custos de terceiros fiquem fora do pacote.

\end{document}
